% ============================================================================
% 中文摘要
% 写作要点：
%   1. 第一段（背景与意义，约130字）：医学问诊实训痛点 + 大模型机遇
%   2. 第二段（系统设计与实现，约180字）：前后端分离架构 + 三角色 +
%      四模块 + 关键基础设施
%   3. 第三段（创新点与结果，约180字）：Agent 编排层 + RAG + Memory +
%      临床工具 + 结构化评分 + 部署测试
%   4. 关键词 6 个，与英文严格对应
%   5. 总字数控制在 500 字左右
% ============================================================================

\begin{abstractzh}

\noindent{\heiti\zihao{-4} 摘要：}问诊训练需要病例、标准化病人与教师反馈共同支撑，但在高校日常教学中，真实病例不易反复开放，标准化病人成本较高，教师也很难逐轮跟进每名学生的训练过程。大语言模型为虚拟病人提供了新的实现途径，不过直接依赖提示词容易出现角色越界、检查结果编造和反馈笼统等问题。基于此，本文围绕医学问诊实训设计并实现虚拟问诊系统，重点解决训练过程可交互、病例事实可约束、训练结果可复盘的问题。

系统采用前后端分离架构。前端基于 Vue 3、TypeScript 和 Element Plus 构建，后端使用 Spring Boot 3.3.6、Java 17 与 MyBatis-Plus 实现；MySQL 保存用户、病例和训练记录，Redis 维护验证码、登录锁定和会话上下文，MinIO 存储医学影像与附件，Docker Compose 用于部署。系统面向学生、教师、管理员三类用户，提供 AI 问诊训练、病例库管理、学习分析和系统管控等功能，并集成 JWT 认证、RBAC 权限控制、ECharts 可视化和 SSE 流式响应。

在 AI 实现上，系统将单一提示词调用拆分为 PatientAgent、DiagnosisAgent 和 FeedbackAgent 三类角色，由 AgentRouter 按训练阶段选择调用对象；关键词检索式 RAG 用于补充医学知识，Redis 滑动窗口保存最近多轮对话，临床工具优先从病例数据返回体格检查、辅助检查和影像结果，避免模型自行生成客观数值。训练结束后，系统以结构化 JSON 输出问诊、诊断、治疗、沟通等维度评分，并在解析失败时使用规则兜底。系统已完成云端部署和功能、性能、安全测试，能够支持医学问诊训练的基本闭环。

\keywordszh{医学教育；虚拟问诊；大语言模型；AI Agent；检索增强生成；前后端分离}

\end{abstractzh}
